\chapter{System Architecture\label{chap:architecture}}

\section{Architecture Overview}

The system follows a layered architecture pattern where each layer has a clearly defined responsibility. The five layers---Simulation, Environment Generation, Task Management, Execution Control, and Formation Following---communicate through well-defined ROS2 interfaces. Figure~\ref{fig:architecture} illustrates the overall system architecture and the relationships between layers.

\begin{figure}[htbp]
  \centering
  \caption{System Architecture Diagram}
  \label{fig:architecture}
  \vspace{0.5em}
  \begin{tabular}{c|c|c}
    \toprule
    \textbf{Layer} & \textbf{Node} & \textbf{Role} \\
    \midrule
    Simulation & \texttt{turtlesim\_node} & Visual simulation environment \\
    Environment Generation & \texttt{spawn\_manager} & Periodic turtle spawning \\
    Task Management & \texttt{master\_manager} & Target selection and dispatch \\
    Execution Control & \texttt{catch\_executor} & Pursuit action execution \\
    Formation Following & \texttt{follower\_manager} & Chain following control \\
    \bottomrule
  \end{tabular}
\end{figure}

The design rationale for this layered decomposition is threefold: (1) each layer has a clear, single responsibility, making the system easy to understand and debug; (2) the communication patterns naturally map to ROS2 paradigms---services for spawning, topics for continuous data flow, and actions for long-running tasks; and (3) the architecture is extensible, allowing individual layers to be modified or replaced without affecting others.

\section{Node Communication Architecture}

The nodes communicate through three ROS2 communication paradigms: \textbf{Topics}, \textbf{Services}, and \textbf{Actions}. Table~\ref{tab:communication} summarises the complete communication interfaces.

\begin{table}[htbp]
\centering
\caption{Node Communication Interfaces}
\label{tab:communication}
\begin{tabular}{llll}
\toprule
\textbf{Interface} & \textbf{Type} & \textbf{Publisher/Client} & \textbf{Subscriber/Server} \\
\midrule
\texttt{/spawn} & Service & \texttt{spawn\_manager} & \texttt{turtlesim\_node} \\
\texttt{/<name>/pose} & Topic & \texttt{turtlesim\_node} & \texttt{master\_manager}, \texttt{catch\_executor}, \texttt{follower\_manager} \\
\texttt{/<name>/cmd\_vel} & Topic & \texttt{catch\_executor}, \texttt{follower\_manager} & \texttt{turtlesim\_node} \\
\texttt{/caught\_chain} & Topic & \texttt{master\_manager} & \texttt{follower\_manager} \\
\texttt{catch\_target} & Action & \texttt{master\_manager} (Client) & \texttt{catch\_executor} (Server) \\
\bottomrule
\end{tabular}
\end{table}

The communication design follows these principles:

\begin{itemize}
  \item \textbf{Topics} are used for continuous, streaming data such as turtle poses and velocity commands, where the latest value is always the most relevant.
  \item \textbf{Services} are used for request--response interactions such as spawning a new turtle, where a confirmation is required.
  \item \textbf{Actions} are used for long-running tasks such as pursuing a target, where real-time feedback and the ability to cancel are essential.
\end{itemize}

\section{Package Structure}

The project is organised into two ROS2 packages:

\begin{itemize}
  \item \texttt{catch\_turtle\_interfaces} (\texttt{ament\_cmake}): Defines the custom action interface \texttt{CatchTarget.action}.
  \item \texttt{catch\_turtle\_bringup} (\texttt{ament\_python}): Contains all Python nodes, launch files, and configuration.
\end{itemize}

The directory structure is as follows:

\begin{lstlisting}[language={},title={Project Directory Structure}]
src/
+-- catch_turtle_interfaces/
|   +-- action/
|   |   +-- CatchTarget.action
|   +-- CMakeLists.txt
|   +-- package.xml
+-- catch_turtle_bringup/
    +-- catch_turtle_bringup/
    |   +-- __init__.py
    |   +-- utils.py
    |   +-- turtle_registry.py
    |   +-- spawn_manager.py
    |   +-- master_manager.py
    |   +-- catch_executor.py
    |   +-- follower_manager.py
    +-- launch/
    |   +-- catch_turtle.launch.py
    +-- config/
    |   +-- params.yaml
    +-- resource/
    |   +-- catch_turtle_bringup
    +-- setup.py
    +-- package.xml
\end{lstlisting}

\section{Concurrency Model}

A critical aspect of the architecture is the concurrency model. Two nodes require multi-threaded execution:

\begin{itemize}
  \item \texttt{catch\_executor} must use \texttt{MultiThreadedExecutor} with \texttt{ReentrantCallbackGroup}, otherwise the long-running Action execute callback would block pose subscriptions, causing the control loop to stall.
  \item \texttt{master\_manager} similarly requires \texttt{MultiThreadedExecutor} to prevent Action callbacks and timer callbacks from blocking each other during preemption operations.
  \item The remaining nodes (\texttt{spawn\_manager} and \texttt{follower\_manager}) operate correctly with single-threaded execution.
\end{itemize}
